\documentclass[11pt,a4paper]{article}
\usepackage[utf8]{inputenc}
\usepackage[T1]{fontenc}
\usepackage{lmodern}
\usepackage[brazilian]{babel}
\usepackage{amsmath,amssymb}
\usepackage{geometry}
\geometry{top=2.0cm, bottom=2.0cm, left=2.0cm, right=2.0cm}
\usepackage{enumitem}
\usepackage{fancyhdr}
\usepackage{titlesec}
\usepackage{xcolor}

\pagestyle{fancy}
\fancyhf{}
\fancyhead[L]{\small\textbf{FSC5101 -- Física I} $\vert$ UFSC}
\fancyhead[R]{\small\textbf{Banco de Questões para Prova} $\vert$ Capítulos 4 e 5}
\fancyfoot[C]{\thepage}
\renewcommand{\headrulewidth}{0.5pt}

\titleformat{\section}{\large\bfseries\color{blue!70!black}}{}{0em}{}[\titlerule]
\titleformat{\subsection}{\bfseries\color{black}}{}{0em}{}

\begin{document}

\begin{center}
    {\LARGE \textbf{UNIVERSIDADE FEDERAL DE SANTA CATARINA}}\\[0.2cm]
    {\large \textbf{Departamento de Física -- FSC5101: Física I}}\\[0.3cm]
    {\Large \textbf{Banco de 40 Questões Selecionadas para Avaliação}}\\[0.2cm]
    {\normalsize \textbf{Fonte:} Sears \& Zemansky -- Física I: Mecânica (Young \& Freedman, 12ª Edição)}\\[0.1cm]
    {\small \textit{Capítulo 4: Leis de Newton do Movimento \quad $\vert$ \quad Capítulo 5: Aplicações das Leis de Newton}}\\[0.2cm]
    \rule{\textwidth}{0.8pt}
\end{center}

\vspace{0.2cm}
\noindent
\textbf{Orientações Gerais:}\\
Este documento contém uma seleção de 40 questões autênticas extraídas diretamente dos Capítulos 4 e 5 do livro-texto \textit{Sears \& Zemansky (Física I, 12ª Ed.)}, divididas por nível de complexidade (10 Fáceis, 15 Médias e 15 Difíceis), destinadas à elaboração e composição de provas. Em todos os cálculos numéricos, adote $g = 9,80\text{ m/s}^2$, exceto quando indicado em contrário no enunciado.

\vspace{0.4cm}

%%%%%%%%%%%%%%%%%%%%%%%%%%%%%%%%%%%%%%%%%%%%%%%%%%%%%%%%%%%%%%%%%%%%%%%%%%%%%%%
\section*{PARTE I: QUESTÕES FÁCEIS (10 QUESTÕES)}
%%%%%%%%%%%%%%%%%%%%%%%%%%%%%%%%%%%%%%%%%%%%%%%%%%%%%%%%%%%%%%%%%%%%%%%%%%%%%%%

\begin{enumerate}[label=\textbf{\arabic*.}, leftmargin=*, itemsep=0.45cm]

% 1
\item \textbf{[Sears 4.7]} Se uma força resultante horizontal de $132\text{ N}$ é aplicada a uma pessoa com massa de $60\text{ kg}$ em repouso na beira de uma piscina, qual é a aceleração produzida?

% 2
\item \textbf{[Sears 4.8]} Qual o módulo da força necessária para imprimir uma aceleração de $1,40\text{ m/s}^2$ em uma geladeira com massa de $135\text{ kg}$?

% 3
\item \textbf{[Sears 4.9]} Uma caixa está em repouso sobre um lago congelado, que é uma superfície horizontal sem atrito. Se um pescador aplica uma força horizontal de módulo $48,0\text{ N}$ sobre a caixa e produz uma aceleração de $3,00\text{ m/s}^2$, qual é a massa da caixa?

% 4
\item \textbf{[Sears 4.10]} Um portuário aplica uma força horizontal constante de $80,0\text{ N}$ a um bloco de gelo sobre uma superfície horizontal lisa. A força de atrito é desprezível. O bloco parte do repouso e se move $11,0\text{ m}$ em $5,0\text{ s}$.
\begin{enumerate}[label=\alph*)]
    \item Qual é a massa do bloco de gelo?
    \item Se o portuário parar de empurrar o bloco depois de $5,0\text{ s}$, qual será a distância percorrida pelo bloco nos $5,0\text{ s}$ posteriores?
\end{enumerate}

% 5
\item \textbf{[Sears 4.19]} Na superfície de Io, uma das luas de Júpiter, a aceleração da gravidade é $g = 1,81\text{ m/s}^2$. Uma melancia pesa $44,0\text{ N}$ na superfície da Terra (onde $g = 9,80\text{ m/s}^2$).
\begin{enumerate}[label=\alph*)]
    \item Qual é a massa da melancia na superfície da Terra?
    \item Qual é a sua massa e o seu peso na superfície de Io?
\end{enumerate}

% 6
\item \textbf{[Sears 4.20]} A mochila de uma astronauta pesa $17,5\text{ N}$ quando ela está na superfície terrestre, mas somente $3,24\text{ N}$ na superfície de um asteroide.
\begin{enumerate}[label=\alph*)]
    \item Qual é a aceleração da gravidade nesse asteroide?
    \item Qual é a massa da mochila no asteroide?
\end{enumerate}

% 7
\item \textbf{[Sears 4.24]} O piso de um elevador exerce uma força normal de $620\text{ N}$ de baixo para cima sobre um passageiro que pesa $650\text{ N}$.
\begin{enumerate}[label=\alph*)]
    \item Quais são as reações dessas duas forças no sentido da Terceira Lei de Newton?
    \item O passageiro está sendo acelerado? Em caso afirmativo, determine o módulo, a direção e o sentido da aceleração.
\end{enumerate}

% 8
\item \textbf{[Sears 5.1]} Dois pesos de $25,0\text{ N}$ estão suspensos nas extremidades opostas de uma corda leve que passa sobre uma polia sem atrito. O centro da polia está ligado a uma corrente presa ao teto.
\begin{enumerate}[label=\alph*)]
    \item Qual é a tensão na corda?
    \item Qual é a tensão na corrente presa ao teto?
\end{enumerate}

% 9
\item \textbf{[Sears 5.11]} Um homem empurra um piano de $180\text{ kg}$, de modo que ele desliza com velocidade constante para baixo de uma rampa inclinada de $11,0^\circ$ acima da horizontal. Despreze o atrito que atua sobre o piano. Calcule o módulo da força aplicada pelo homem se ela for:
\begin{enumerate}[label=\alph*)]
    \item paralela ao plano inclinado;
    \item paralela ao piso horizontal.
\end{enumerate}

% 10
\item \textbf{[Sears 5.28]} Uma caixa de $25,0\text{ kg}$ com livros repousa sobre uma mesa horizontal. O coeficiente de atrito estático entre a caixa e a mesa é $0,40$ e o coeficiente de atrito cinético é $0,20$.
\begin{enumerate}[label=\alph*)]
    \item Qual é o peso da caixa e o módulo da força normal exercida pela mesa sobre ela?
    \item Qual força horizontal é necessária para iniciar o movimento da caixa?
    \item Qual força horizontal é necessária para manter a caixa em movimento com velocidade constante uma vez iniciado o deslizamento?
\end{enumerate}

\end{enumerate}

\newpage
%%%%%%%%%%%%%%%%%%%%%%%%%%%%%%%%%%%%%%%%%%%%%%%%%%%%%%%%%%%%%%%%%%%%%%%%%%%%%%%
\section*{PARTE II: QUESTÕES MÉDIAS (15 QUESTÕES)}
%%%%%%%%%%%%%%%%%%%%%%%%%%%%%%%%%%%%%%%%%%%%%%%%%%%%%%%%%%%%%%%%%%%%%%%%%%%%%%%

\begin{enumerate}[label=\textbf{\arabic*.}, leftmargin=*, itemsep=0.45cm, resume]

% 11
\item \textbf{[Sears 4.5]} Dois cachorros puxam horizontalmente cordas amarradas a um poste; o ângulo entre as duas cordas é igual a $60,0^\circ$. Se o cachorro A exerce uma força de $270\text{ N}$ e o cachorro B exerce uma força de $300\text{ N}$, ache o módulo da força resultante sobre o poste e o ângulo que ela forma com a corda puxada pelo cachorro A.

% 12
\item \textbf{[Sears 4.11]} Um disco de hóquei com massa de $0,160\text{ kg}$ está em repouso na origem ($x = 0$) em uma superfície horizontal sem atrito da pista. No instante $t = 0$, um jogador aplica uma força constante de $0,250\text{ N}$ sobre o disco, paralela ao eixo $Ox$. Ele continua a aplicar essa força até o instante $t = 2,00\text{ s}$.
\begin{enumerate}[label=\alph*)]
    \item Qual é a posição e a velocidade do disco no instante $t = 2,00\text{ s}$?
    \item Se a mesma força for aplicada novamente no instante $t = 5,00\text{ s}$ no mesmo sentido, qual será a posição e a velocidade do disco no instante $t = 7,00\text{ s}$?
\end{enumerate}

% 13
\item \textbf{[Sears 4.17]} Uma caixa com massa de $5,0\text{ kg}$ é puxada por uma corda ao longo de uma superfície horizontal lisa sem atrito. A tensão na corda é de $40,0\text{ N}$ e ela faz um ângulo de $30,0^\circ$ acima da horizontal.
\begin{enumerate}[label=\alph*)]
    \item Qual é a aceleração horizontal da caixa?
    \item Qual é o módulo da força normal que a superfície exerce sobre a caixa?
\end{enumerate}

% 14
\item \textbf{[Sears 4.25]} Uma estudante com massa de $45\text{ kg}$ pula de um trampolim elevado. Considerando a massa da Terra como $6,0 \times 10^{24}\text{ kg}$, qual é a aceleração da Terra no sentido da estudante quando ela se acelera no sentido da Terra com $9,80\text{ m/s}^2$? Suponha que a força resultante sobre a Terra seja a força gravitacional que a estudante exerce sobre ela.

% 15
\item \textbf{[Sears 4.36]} Um estudante de física com $550\text{ N}$ de peso sobe em uma balança de mola apoiada no chão de um elevador. Conforme o elevador se move, a balança registra uma leitura de peso aparente de $450\text{ N}$.
\begin{enumerate}[label=\alph*)]
    \item Determine o módulo, a direção e o sentido da aceleração do elevador.
    \item Se a balança registrar $670\text{ N}$, qual será a aceleração do elevador?
    \item Se a leitura da balança for nula, o que se pode concluir sobre o movimento do elevador?
\end{enumerate}

% 16
\item \textbf{[Sears 4.39]} Dois blocos de massas $m_1 = 4,0\text{ kg}$ e $m_2 = 6,0\text{ kg}$ estão apoiados lado a lado em contato mútuo sobre uma superfície horizontal sem atrito. Uma força horizontal constante $F = 30,0\text{ N}$ é aplicada ao bloco $m_1$, empurrando-o contra o bloco $m_2$.
\begin{enumerate}[label=\alph*)]
    \item Desenhe o diagrama do corpo livre para cada um dos blocos.
    \item Qual é a aceleração do sistema formado pelos dois blocos?
    \item Qual é o módulo da força de contato exercida entre os dois blocos?
\end{enumerate}

% 17
\item \textbf{[Sears 5.3]} Um quadro está suspenso em uma parede por meio de dois fios presos aos seus cantos superiores. Cada fio forma um ângulo de $30,0^\circ$ com a vertical. Se a massa total do quadro for de $1,50\text{ kg}$, calcule a tensão em cada um dos fios.

% 18
\item \textbf{[Sears 5.9]} Um carro de $1000\text{ kg}$ está sobre uma rampa sem atrito inclinada de $20,0^\circ$ em relação à horizontal, sendo sustentado por um cabo de aço paralelo à rampa.
\begin{enumerate}[label=\alph*)]
    \item Se o carro sobe a rampa com velocidade constante de $15,0\text{ m/s}$, qual é a tensão no cabo?
    \item Se o carro sobe a rampa acelerando a uma taxa de $2,00\text{ m/s}^2$, qual é a tensão no cabo?
\end{enumerate}

% 19
\item \textbf{[Sears 5.17]} Um bloco de gelo de $8,00\text{ kg}$, liberado a partir do repouso no topo de uma rampa sem atrito com $1,50\text{ m}$ de comprimento, desliza para baixo, alcançando uma velocidade escalar de $2,50\text{ m/s}$ na base da rampa.
\begin{enumerate}[label=\alph*)]
    \item Qual é o ângulo de inclinação da rampa com a horizontal?
    \item Qual seria a velocidade escalar do bloco na base se ele fosse lançado para cima a partir da base ao longo da mesma rampa com velocidade inicial de $2,50\text{ m/s}$?
\end{enumerate}

% 20
\item \textbf{[Sears 5.19]} \textbf{Máquina de Atwood.} Uma carga de tijolos com massa de $15,0\text{ kg}$ é suspensa pela extremidade de uma corda leve que passa sobre uma pequena polia fixa sem atrito. Um contrapeso de $28,0\text{ kg}$ está preso na outra extremidade da corda. O sistema é liberado a partir do repouso.
\begin{enumerate}[label=\alph*)]
    \item Desenhe o diagrama do corpo livre para a carga de tijolos e outro para o contrapeso.
    \item Qual é o módulo da aceleração do sistema?
    \item Qual é a tensão na corda durante o movimento acelerado?
\end{enumerate}

% 21
\item \textbf{[Sears 5.23]} Um tobogã com massa de $75,0\text{ kg}$ desce uma encosta coberta de neve inclinada de $25,0^\circ$ com a horizontal. O coeficiente de atrito cinético entre os patins do tobogã e a neve é $\mu_c = 0,12$.
\begin{enumerate}[label=\alph*)]
    \item Desenhe o diagrama do corpo livre para o tobogã.
    \item Calcule o módulo da força normal exercida pela encosta sobre o tobogã.
    \item Calcule a aceleração do tobogã ao longo da encosta.
\end{enumerate}

% 22
\item \textbf{[Sears 5.31]} Uma caixa de ferramentas de $45,0\text{ kg}$ está em repouso sobre um piso horizontal. Você exerce sobre ela uma força de puxar horizontal que aumenta gradualmente e observa que a caixa só começa a se mover quando a sua força ultrapassa $313\text{ N}$. A partir daí, você deve reduzir sua força de puxar para $208\text{ N}$ para mantê-la em movimento com velocidade constante de $25,0\text{ cm/s}$.
\begin{enumerate}[label=\alph*)]
    \item Quais são os coeficientes de atrito estático ($\mu_e$) e cinético ($\mu_c$) entre a caixa e o piso?
    \item Que força horizontal de puxar você deve exercer para imprimir à caixa uma aceleração de $1,10\text{ m/s}^2$?
    \item Suponha que você realize essa experiência na Lua ($g = 1,62\text{ m/s}^2$). Qual seria a força necessária para iniciar o movimento?
\end{enumerate}

% 23
\item \textbf{[Sears 5.33]} Um engradado de $48,0\text{ kg}$ está sobre um piso horizontal rugoso com coeficiente de atrito cinético $\mu_c = 0,30$. Você puxa o engradado por meio de uma corda que faz um ângulo de $30,0^\circ$ acima da horizontal. Se a tensão na corda for de $250\text{ N}$:
\begin{enumerate}[label=\alph*)]
    \item Calcule o módulo da força normal exercida pelo piso sobre o engradado.
    \item Calcule a força de atrito cinético exercida sobre o engradado.
    \item Calcule a aceleração horizontal do engradado.
\end{enumerate}

% 24
\item \textbf{[Sears 5.43]} Uma curva plana em uma rodovia possui raio de curvatura igual a $220\text{ m}$. Em um dia chuvoso, o coeficiente de atrito estático entre os pneus de um veículo e o asfalto é $\mu_e = 0,40$. Qual é a velocidade máxima com que um automóvel pode realizar essa curva sem derrapar lateralmente (expresse a resposta em m/s e em km/h)?

% 25
\item \textbf{[Sears 5.47]} \textbf{Pêndulo Cônico.} Uma pequena esfera pesada de massa $m$ está presa à extremidade de uma corda leve de comprimento $L = 1,20\text{ m}$. A esfera gira com velocidade constante em uma trajetória circular horizontal de raio $r = 0,50\text{ m}$, de modo que a corda descreve a superfície de um cone reto formando um ângulo $\beta$ com a vertical.
\begin{enumerate}[label=\alph*)]
    \item Determine o valor do ângulo $\beta$.
    \item Calcule a velocidade escalar constante da esfera ao longo do círculo.
    \item Determine o período de revolução $T$ do movimento circular.
\end{enumerate}

\end{enumerate}

\newpage
%%%%%%%%%%%%%%%%%%%%%%%%%%%%%%%%%%%%%%%%%%%%%%%%%%%%%%%%%%%%%%%%%%%%%%%%%%%%%%%
\section*{PARTE III: QUESTÕES DIFÍCEIS -- PROBLEMAS (15 QUESTÕES)}
%%%%%%%%%%%%%%%%%%%%%%%%%%%%%%%%%%%%%%%%%%%%%%%%%%%%%%%%%%%%%%%%%%%%%%%%%%%%%%%

\begin{enumerate}[label=\textbf{\arabic*.}, leftmargin=*, itemsep=0.45cm, resume]

% 26
\item \textbf{[Sears 4.45]} Para estudar o dano que a colisão com grandes pássaros pode causar a um avião, projeta-se uma arma de teste pneumática que acelera projéteis ao longo de um cano cilíndrico reto, de modo que a posição do projétil em função do tempo é dada por:
\[ x(t) = (9,0 \times 10^3\text{ m/s}^2)\,t^2 - (8,0 \times 10^4\text{ m/s}^3)\,t^3 \]
O projétil de massa $m = 1,50\text{ kg}$ deixa a boca do cano no instante $t = 0,025\text{ s}$.
\begin{enumerate}[label=\alph*)]
    \item Qual deve ser o comprimento do cano da arma de teste?
    \item Qual é a velocidade escalar do projétil no instante em que ele emerge do cano?
    \item Qual é o módulo da força resultante que atua sobre o projétil em: i) $t = 0$; ii) $t = 0,025\text{ s}$?
\end{enumerate}

% 27
\item \textbf{[Sears 4.47]} Um instrumento científico de $6,50\text{ kg}$ está pendurado por um cabo vertical no interior de uma espaçonave de teste que está sendo lançada verticalmente a partir da superfície terrestre. O foguete parte do repouso e atinge a altitude de $276\text{ m}$ em $15,0\text{ s}$ com aceleração vertical constante.
\begin{enumerate}[label=\alph*)]
    \item Desenhe o diagrama do corpo livre para o instrumento durante essa fase do voo, indicando claramente qual força possui maior módulo.
    \item Calcule o módulo da força de tração que o cabo vertical exerce sobre o instrumento.
\end{enumerate}

% 28
\item \textbf{[Sears 4.53]} Um paraquedista com massa total de $75,0\text{ kg}$ pula de uma aeronave a grande altitude. Durante a queda vertical antes de abrir o paraquedas, a força de resistência do ar oposta à velocidade é dada por $F_{\text{res}} = -k v^2$, onde $k = 0,250\text{ kg/m}$.
\begin{enumerate}[label=\alph*)]
    \item Obtenha o valor da velocidade terminal $v_t$ atingida pelo paraquedista.
    \item Determine a aceleração instantânea de queda quando a sua velocidade for igual a metade da velocidade terminal ($v = v_t / 2$).
\end{enumerate}

% 29
\item \textbf{[Sears 4.59]} Um objeto com massa $m$ move-se em linha reta ao longo do eixo $Ox$. Sua coordenada de posição em função do tempo é descrita por $x(t) = A t - B t^3$, onde $A$ e $B$ são constantes positivas conhecidas.
\begin{enumerate}[label=\alph*)]
    \item Obtenha as expressões algébricas para a velocidade $v(t)$ e a aceleração $a(t)$ do objeto.
    \item Encontre a expressão analítica da força resultante $F(t)$ exercida sobre o corpo.
    \item Em qual instante de tempo a força resultante sobre o corpo se anula?
\end{enumerate}

% 30
\item \textbf{[Sears 5.57]} Um trabalhador de $80,0\text{ kg}$ está de pé sobre uma plataforma suspensa de $25,0\text{ kg}$ e puxa uma corda leve que passa por uma polia presa ao teto, de modo que ele e a plataforma aceleram juntos para cima com taxa constante de $1,20\text{ m/s}^2$. Despreze o atrito e as massas das polias e da corda.
\begin{enumerate}[label=\alph*)]
    \item Desenhe os diagramas do corpo livre separados para o homem e para a plataforma.
    \item Calcule a força de tração que o homem deve exercer sobre a corda.
    \item Calcule o módulo da força normal que atua entre os pés do homem e a plataforma.
\end{enumerate}

% 31
\item \textbf{[Sears 5.63]} \textbf{Corda com massa.} Em problemas de mecânica, as massas dos cabos e cordas costumam ser desprezadas. Quando uma corda pesada e flexível de massa $M$ está suspensa horizontalmente entre dois apoios fixados à mesma altura, suas extremidades formam um ângulo $\theta$ com a linha horizontal.
\begin{enumerate}[label=\alph*)]
    \item Determine a tensão $T_1$ nas extremidades da corda em função de $M$, $g$ e $\theta$.
    \item Determine a tensão $T_0$ no ponto mais baixo (ponto médio inferior) da corda.
    \item Explique por que é fisicamente impossível que o ângulo $\theta$ seja rigorosamente igual a zero para uma corda com massa sob a ação da gravidade.
\end{enumerate}

% 32
\item \textbf{[Sears 5.70]} Um foguete de pesquisa de $25\,000\text{ kg}$ é lançado verticalmente da superfície terrestre com aceleração constante dirigida para cima de $2,50\text{ m/s}^2$. Durante a fase inicial do movimento, adote $g = 9,80\text{ m/s}^2$.
\begin{enumerate}[label=\alph*)]
    \item Qual é o módulo da força total de propulsão exercida pelos motores sobre o foguete?
    \item Um sensor experimental de $15,0\text{ kg}$ está apoiado no assoalho da câmara de carga do foguete. Calcule a força normal que o piso exerce sobre o sensor durante essa subida acelerada.
\end{enumerate}

% 33
\item \textbf{[Sears 5.71]} Dois blocos retangulares de massas $m_A = 2,0\text{ kg}$ e $m_B = 4,0\text{ kg}$ estão sobrepostos: o bloco A está apoiado diretamente sobre o bloco B. O coeficiente de atrito estático entre as superfícies em contato dos blocos A e B é $\mu_e = 0,30$. A superfície da mesa horizontal sobre a qual o bloco B se apoia é perfeitamente lisa (sem atrito). Uma força horizontal $F$ é aplicada diretamente sobre o bloco inferior B.
\begin{enumerate}[label=\alph*)]
    \item Qual é a aceleração máxima que o bloco A pode ter antes de começar a escorregar sobre o bloco B?
    \item Qual é o módulo máximo da força horizontal $F$ que pode ser aplicada ao bloco B para que os blocos A e B se acelerem juntos sem que haja deslizamento entre eles?
\end{enumerate}

% 34
\item \textbf{[Sears 5.77]} Um bloco com massa $m = 3,00\text{ kg}$ está sobre uma rampa inclinada de $37,0^\circ$ com a horizontal. O coeficiente de atrito cinético entre o bloco e a rampa é $\mu_c = 0,30$. O bloco é lançado para cima ao longo da rampa com velocidade inicial $v_0 = 8,00\text{ m/s}$.
\begin{enumerate}[label=\alph*)]
    \item Calcule a distância que o bloco percorre rampa acima até parar momentaneamente.
    \item Verifique se o bloco, após parar no ponto mais alto, permanece em repouso ou desliza de volta para baixo, sabendo que $\mu_e = 0,40$. Em caso afirmativo, determine sua velocidade ao retornar à base.
\end{enumerate}

% 35
\item \textbf{[Sears 5.85]} Três blocos estão conectados por meio de fios leves e polias ideais sem atrito: o bloco A de massa $m_A = 4,0\text{ kg}$ está suspenso verticalmente na extremidade esquerda; o bloco B de massa $m_B = 12,0\text{ kg}$ desliza sobre uma mesa horizontal com coeficiente de atrito cinético $\mu_c = 0,25$; e o bloco C está suspenso verticalmente na extremidade direita. O bloco B está se movendo para a direita com uma aceleração constante de $2,00\text{ m/s}^2$.
\begin{enumerate}[label=\alph*)]
    \item Determine a tensão na corda que conecta o bloco A ao bloco B.
    \item Determine a massa necessária para o bloco pendente C.
\end{enumerate}

% 36
\item \textbf{[Sears 5.89]} Dois objetos com massas de $5,0\text{ kg}$ e $2,0\text{ kg}$ estão suspensos a $0,600\text{ m}$ acima do solo, presos nas extremidades opostas de uma corda leve de $6,0\text{ m}$ que passa sobre uma polia fixa sem atrito. Os dois corpos partem do repouso. O objeto mais pesado desce até colidir com o solo e permanecer em repouso. Calcule a altura máxima que o objeto de $2,0\text{ kg}$ atinge acima do solo.

% 37
\item \textbf{[Sears 5.91]} Um bloco A com massa $m$ é pressionado contra a face vertical dianteira de um carrinho de transporte. O coeficiente de atrito estático entre o bloco e a face vertical do carrinho é $\mu_e$.
\begin{enumerate}[label=\alph*)]
    \item Qual deve ser a aceleração mínima do carrinho para que o bloco A não escorregue para baixo em relação à parede vertical do carrinho?
    \item Descreva como as forças sobre o bloco são justificadas e compreendidas por um passageiro situado em repouso dentro do referencial do carrinho acelerado.
\end{enumerate}

% 38
\item \textbf{[Sears 5.95]} \textbf{Curva inclinada com atrito.} Uma curva com raio de curvatura igual a $120\text{ m}$ em uma estrada possui uma sobrelevação (ângulo de inclinação lateral $\theta$) projetada para que veículos a uma velocidade escalar de $20,0\text{ m/s}$ contornem a curva sem depender de nenhuma força de atrito lateral.
\begin{enumerate}[label=\alph*)]
    \item Calcule o ângulo de sobrelevação lateral $\theta$ da rodovia.
    \item Se um veículo contornar essa mesma curva com velocidade de $30,0\text{ m/s}$, qual deve ser o coeficiente de atrito estático mínimo ($\mu_e$) entre os pneus e a pista para que o carro não derrape para fora da curva?
\end{enumerate}

% 39
\item \textbf{[Sears 5.106]} \textbf{Resistência de fluidos.} Uma pequena esfera pesada move-se no interior de um fluido viscoso, sofrendo uma força de resistência proporcional à velocidade escalar instantânea: $F_{\text{res}} = -k v$. Sabe-se que a velocidade terminal de descida da esfera é $v_t = 2,00\text{ m/s}$. A esfera é projetada verticalmente de baixo para cima dentro do fluido a partir da origem com velocidade inicial $v_0 = 6,00\text{ m/s}$. Despreze o empuxo.
\begin{enumerate}[label=\alph*)]
    \item Na ausência de resistência do fluido, qual seria a altura máxima atingida pela esfera e quanto tempo ela levaria para atingir o topo?
    \item Levando em conta a resistência proporcional do fluido, deduza a equação da velocidade em função do tempo $v(t)$ e calcule a altura máxima atingida pela esfera no fluido.
\end{enumerate}

% 40
\item \textbf{[Sears 5.121]} Uma cunha triangular de massa $M$ e ângulo de inclinação $\alpha$ repousa sobre uma superfície horizontal perfeitamente lisa sem atrito. Um pequeno bloco de massa $m$ é colocado sobre a rampa da cunha, também sem atrito entre o bloco e a cunha. Uma força horizontal constante $F$ é aplicada à cunha. Qual deve ser o módulo da força $F$ para que o bloco permaneça em repouso em relação à cunha (isto é, mantendo-se a uma altura constante em relação à base e sem deslizar rampa acima nem rampa abaixo)?

\end{enumerate}

\end{document}
